\uuid{o45G}
\titre{Injection/surjection}
\auteur{bodin ; hanani ; mzari}
\organisation{exo7}
\licence{CC BY-NC-SA 4.0}
\datecreate{2020-01-20}
\niveau{L1}
\module{Algèbre}
\chapitre{Logique, ensemble, raisonnement}
\sousChapitre{Ensemble}
\difficulte{2}
\theme{Ensembles, applications, injection, surjection}
\exo7qcmid{lille-1.ensembles.020}

\contenu{
\begin{question}
\qclassification{100.02,101.01,102.02,102.02}
\qtags{motcle=injection/surjection}
On considère l'application $f:\Rr\setminus\{2\}\to \Rr\setminus \{1\}$ définie par
$$\forall x\in \Rr\setminus\{2\},\; f(x)=\frac{x+1}{x-2}.$$
Quelle est la bonne réponse ?
\begin{answers}
    \bad{$f$ n'est pas bijective.}
    \bad{$f$ est bijective et $\displaystyle f^{-1}(x)=\frac{x-2}{x+1}$.}
    \good{$f$ est bijective et $\displaystyle f^{-1}(x)=\frac{2x+1}{x-1}$.}
    \bad{$f$ est bijective et $\displaystyle f^{-1}(x)=\frac{-x+1}{-x-2}$.}
\end{answers}
\begin{explanations}
Tout $y\neq 1$ admet un unique antécédent qui s'écrit $\displaystyle x=\frac{2y+1}{y-1}\in \Rr\setminus\{2\}$. Donc $f$ est bijective et $\displaystyle f^{-1}(y)=\frac{2y+1}{y-1}$.
\end{explanations}
\end{question}
}
